%% SafeTransition — EMS software description manuscript (paper 3), v1
%% Companion: paper 1 "Aggregate Indices and Transition Safety" (v46),
%%            paper 2 "An Obstruction Calculus for Viability under Incomplete Observation" (v43).
%% Compile: tectonic paper3_safetransition_ems_v1.tex  (figures under ../figs_p3 resolved via \graphicspath)
\documentclass[3p,11pt]{elsarticle}
\usepackage{amsmath,amssymb}
\usepackage{booktabs}
\usepackage{graphicx}
\usepackage{hyperref}
\hypersetup{colorlinks=true, urlcolor=blue, citecolor=blue, linkcolor=blue}
\emergencystretch=2.5em

\graphicspath{{/home/user/figs_p3/}}

\journal{Environmental Modelling \& Software}

\begin{document}

\begin{frontmatter}

\title{SafeTransition: exact rational certification of transition safety for
sustainability assessment}

\author[aff]{Amin Abaee}
\address[aff]{Independent Researcher,
\href{https://orcid.org/0000-0002-0019-1842}{ORCID: 0000-0002-0019-1842},
\href{mailto:amin\_abaee@ut.ac.ir}{amin\_abaee@ut.ac.ir}}

\begin{abstract}
SafeTransition, a Python library for certifying the transition safety of
sustainability assessments, is presented. The library implements the
assessment-operator framework of the companion manuscripts: typed floor
constraints evaluated on exact reachability tubes, scalarized aggregate
operators at arbitrary positive weights, and backward recursions over state
and belief spaces. Infeasibility is returned as a Farkas certificate, a
nonnegative multiplier vector with vanishing weighted constraint sum and
negative weighted bound, re-verified exactly together with its infeasibility
margin. All computations run in exact rational arithmetic with no runtime
dependencies beyond the Python standard library; floating point appears only
in rendered graphics. Twenty-four exact checks re-derive the
resource-transition benchmark of the companion verification deposit,
including the per-weight licensing thresholds $2/3$ and $3/2$ and the rescue
threshold $1-x$. A single-file dashboard renders the resulting readings.
The library is distributed under the MIT licence.
\end{abstract}

\begin{keyword}
sustainability assessment \sep viability theory \sep transition safety \sep
Farkas certificate \sep exact arithmetic \sep composite indicators
\sep multi-criteria decision analysis
\end{keyword}

\end{frontmatter}

\section{Introduction}\label{sec:introduction}

Composite indicators occupy a central place in sustainability monitoring.
Aggregates reduce heterogeneous floor constraints on income, biomass, and
social outcomes to a single weighted reading, and that reading is what
dashboards, statutory reviews, and public communications display. The
practice rests on a long modelling tradition, and its governance follows
established frameworks for good model development and evaluation practice
\citep{Jakeman2006}. It also rests on a mathematical step that
is easy to audit and easy to overlook: a weighted sum is certified positive
at every audited snapshot, and the system behind the snapshots is declared
safe.

The companion manuscripts show that this step can fail, and characterize
exactly when it fails. In the assessment-operator framework of \citet{Abaee2026a},
a noncompensatory typed operator --- which requires every floor to hold on
the full trajectory, not at endpoints alone --- can reject every plan that a
scalarized aggregate operator licenses at a given weight; the separation is
exhibited on an explicit rational witness datum and propagated through
backward recursions. A resource-transition benchmark realizes the witness as
a regulated fishery transition and verifies every quoted number in exact
rational arithmetic. Under incomplete observation, the obstruction calculus
of \citet{Abaee2026b} complements viability assessment from the other side: it
certifies that no observation-based policy is viable, through finitely
checkable certificates --- the common-action obstruction with an explicit
Farkas multiplier vector, the observation-fibre criterion, and a
finite-horizon belief-space recursion that is sound and complete in finite
systems. These results are mathematical; they describe what should be
computed. They do not yet describe a tool that computes it.

This is the gap SafeTransition addresses. Reading a composite dashboard
today, an analyst can recompute the aggregate with ordinary floating-point
tooling, but has no standard way to re-derive the safety verdict behind it,
to check that an infeasible plan menu is infeasible by certificate rather
than by solver exit code, or to see which weight ratios license which plans.
Numerical viability and control toolboxes compute kernels and invariant sets
to high precision; multi-criteria decision analysis packages aggregate and
rank; formal verification systems certify finite-state models against
specifications. What is missing in the intersection is a small, auditable
layer in which the verdicts themselves --- admissibility of a plan at a
state, licensing at a weight, emptiness of a menu, emptiness of an epistemic
kernel --- are returned as exact rational objects that a third party can
re-verify by hand or by machine.

This paper presents SafeTransition, a Python library that provides this
layer. The library implements the five assessment operators and their
accepted-state sets, the typed finite-graph backward recursion, the
belief-space robust epistemic recursion, and the finitely checkable
certificate family: exact rational Fourier--Motzkin elimination with
provenance tracking, returning Farkas certificates and infeasibility
margins; the finite common-action obstruction with a minimal conflicting
subfamily; and the observation-fibre criterion. Every computation runs in
exact rational arithmetic; floats appear only in rendered graphics. The
library re-derives, in twenty-four exact checks, the resource-transition
benchmark of the verification deposit, and renders the resulting readings as
a single-file dashboard. It depends on nothing beyond the Python standard
library, which makes the exactness claim structural rather than procedural:
there is no numerical backend whose rounds of approximation would need to be
audited.

The paper is a software description. Section 2 states the mathematical basis
as implemented, with proofs delegated to the companion manuscripts.
Section 3 describes the software architecture and the exact-arithmetic
discipline. Section 4 illustrates the library on the benchmark, on the
dashboard readings, and on the partial-observation demonstrations. Section 5
positions the library relative to existing software, and Section 6 records
availability and requirements. Section 7 concludes.

\section{Mathematical basis}\label{sec:math}

\subsection{Transition data and assessment operators}\label{sec:operators}

A typed transition datum is a phase state $z = (q, x, s_1, \dots, s_n)$ with
a phase flag $q$, fund coordinates $x$, and typed floors $s_i$ normalized to
zero; a finite action set; a set of disturbances with, for each action and
disturbance, an exact tube $\mathrm{Tube}(a, d)$ --- the set of states
visited over the review period, given here as piecewise-linear coordinate
paths on breakpoints at which every visited extreme is attained; a
transition-safe set $S$; and a destination set $G$ reached under an endpoint
reset. Actions in the library declare offsets on the breakpoints, together
with an action-indexed worst-case disturbance, so that each tube is the
exact visited set and no outer approximation enters.

Five operators are distinguished, following \citet{Abaee2026a}. For a state $z$,
the noncompensatory typed operator admits action $a$ when every floor holds
on $\mathrm{Tube}(a, d)$ for every disturbance $d$ and the successor lies in
$G$; the scalarized aggregate operator replaces the coordinate-wise
condition by $w \cdot s \geq 0$ at weight $w$ in the positive weight cone;
the exact-tube physical and endpoint-only physical operators weaken the
constraint set; and the typed-endpoint operator evaluates floors at
endpoints only. Because endpoints are contained in tubes and typed sets in
physical sets, the chain
\[
E_{\mathrm{typ}}(z) \;\subseteq\; E_w(z) \;\subseteq\;
E_{\mathrm{tube,phys}}(z) \;\subseteq\; E_{\mathrm{end}}(z)
\]
holds for every positive weight, and the library verifies it exactly at any
state on request (Section~\ref{sec:api}). The accepted-state set
$\mathcal{V}[E] = \{ z : E(z) \neq \emptyset \}$ and the compensatory
accepted set $\mathcal{V}_{\mathrm{weak}} = \bigcap_w \mathcal{V}_w$ are
derived from the operators; the witnessed separation
$\mathcal{V}_{\mathrm{typ}} \subseteq \mathcal{V}_{\mathrm{weak}}$, strict in
general, is what a composite dashboard cannot see (Section~\ref{sec:benchmark}).

\subsection{Recursions}\label{sec:recursions}

Two backward recursions are implemented. The typed recursion propagates
viability backwards through an explicit state--action graph: a state is
$k$-step viable when some action keeps every visited state inside the safe
set and lands in the $(k-1)$-step viable set. On explicit graphs with
constant-time predicates the recursion is polynomial in graph size and
horizon. The belief-space recursion addresses partially observed systems,
following the robust epistemic construction of \citet{Abaee2026b}: states carry
observation labels, an action is admissible for a belief when no state in
the belief's fibre enters a violation en route, and beliefs are updated to
the labels reachable under the action. The recursion returns the sets
$\mathcal{W}_k$ of $k$-step viable beliefs and is sound and complete in
finite systems: a belief outside $\mathcal{W}_N$ admits no viable
observation-based policy over the horizon. Post-observation recourse
failures --- beliefs in $\mathcal{W}_1 \setminus \mathcal{W}_2$ --- are the
signature output (Section~\ref{sec:partial}).

\subsection{Certificates}\label{sec:certificates}

Infeasibility is certified, not merely reported. For a stacked rational
system $A u \leq b$, the library runs Fourier--Motzkin elimination with
provenance tracking: every derived row carries its multiplier vector over
the original rows, and when a contradiction $0 \leq c$, $c < 0$ is derived,
the normalized multipliers $\lambda \geq 0$ with $\lambda^{\top} A = 0$ and
$\lambda^{\top} b < 0$ are returned as a Farkas certificate, together with
the infeasibility margin $-\lambda^{\top} b > 0$. The certificate is
re-verified exactly before being returned. The finite common-action
obstruction fires when the safe-action sets of pairwise-compatible states
intersect emptily while each is nonempty --- every state is individually
viable under full information, yet no observation-based policy is viable;
the library reports a minimal conflicting subfamily. The observation-fibre
criterion decides whether an exact observation-only certifier can exist at
all: one exists if and only if safe-set membership is constant on
observation fibres, and a violating fibre is reported otherwise.

\section{Software description}\label{sec:software}

\subsection{Architecture}\label{sec:architecture}

The library is organized as nine modules (Table~\ref{tab:modules}). The
\texttt{datum} module defines the transition-data interface and ships one
fully specified instance, the witness datum of the companion manuscript,
with destination gain $e = (1/4, 1/4)$ and rescue cost $c = 1$. Custom data
subclass the interface. The \texttt{operators}, \texttt{recursion}, and
\texttt{certificates} modules implement Section~\ref{sec:math}; the
\texttt{indicators} and \texttt{dashboard} modules produce the readings of
Section~\ref{sec:readings}; the \texttt{benchmark} module carries the
twenty-four exact checks of Section~\ref{sec:benchmark}. A command-line
interface exposes \texttt{verify}, \texttt{report}, and \texttt{demo}.

\begin{table}[t]
\centering
\small
\begin{tabular}{@{}ll@{}}
\toprule
Module & Contents \\
\midrule
\texttt{rational} & exact rational helpers (interpretation, rendering) \\
\texttt{datum} & typed transition data; the witness datum instance \\
\texttt{operators} & five operators, accepted-state sets, chain check \\
\texttt{recursion} & typed finite-graph recursion; belief-space recursion \\
\texttt{certificates} & Fourier--Motzkin with provenance; Farkas verification; \\
                       & common-action obstruction; fibre criterion \\
\texttt{indicators} & licensing thresholds, rescue threshold, blindness alarm \\
\texttt{dashboard} & single-file HTML rendering (inline SVG/CSS) \\
\texttt{benchmark} & twenty-four exact checks; verified schedule values \\
\texttt{cli} & \texttt{verify} / \texttt{report} / \texttt{demo} \\
\bottomrule
\end{tabular}
\caption{Module organization of SafeTransition 1.0.0
(1\,269 source lines, 283 test lines).}
\label{tab:modules}
\end{table}

\subsection{The exact-arithmetic discipline}\label{sec:exactness}

Every assessment, recursion, and certificate is computed in exact rational
arithmetic. Rationals enter either as integers, as decimal strings, or as
explicit numerator--denominator pairs; no implicit conversion from binary
floating point exists anywhere in the library, so an input's exactness is
structural, not a matter of reviewer discipline. Floats occur in exactly two
places, both rendering: SVG geometry in the dashboard and matplotlib
coordinates in figure scripts. Checks never read floats. The distinction is
the library's design-centre claim: outputs are verifiable, not merely
computed. A Farkas certificate, a licensing threshold, or a benchmark check
can be re-verified by any rational-arithmetic implementation --- including
pencil and paper --- and the test suite exercises exactly this: every
certificate returned by the eliminator is re-verified against its defining
identities before being returned, and the benchmark checks re-derive each
deposited value symbolically (for example $\sigma(16/5) = 1088/125$ for the
sustained-yield quota of the Schaefer realization).

\subsection{Interfaces and complexity}\label{sec:api}

The application programming interface mirrors the mathematics:
\texttt{E(datum, z, mode, w)} returns the action set of an operator at a
state; \texttt{check\_chain} verifies the operator chain exactly;
\texttt{belief\_backward} returns the belief-viability sets
$\mathcal{W}_k$; \texttt{certify\_polyhedron} returns a feasibility verdict
with its certificate; \texttt{compute} returns the dashboard readings at a
state, weight, and read plan. The typed recursion on explicit graphs is
polynomial in graph size and horizon; the Fourier--Motzkin eliminator is
exponential in the worst case, as any exact general method is, which is why
the certificate family targets the small stacked systems that transition
menus produce (a handful of constraints; the witness instances eliminate in
milliseconds). Belief enumeration bounds itself (4\,096 reachable beliefs)
and raises rather than drifting.

\section{Illustrative examples}\label{sec:examples}

\subsection{The resource-transition benchmark}\label{sec:benchmark}

The benchmark realizes the witness datum as a regulated fishery transition:
stock biomass with Schaefer surplus $\sigma(B) = rB(1 - B/K)$ ($r = 4$,
$K = 10$, $B_{\mathrm{lim}} = 2$), a pulse-plus-closed-season quota plan,
a mid-season heatwave strike, and a staged rebuild financed by a fund
buy-back. Running \texttt{safetransition verify} executes twenty-four exact
checks that re-derive the deposited verification values: the witness datum
and tube tables; quota admissibility on both disturbance branches
($0 \leq H^* \leq 13$, peak $1463/125$); exact tracking of the benign
branch and conservatism of the certified tubes for the nonlinear
realization ($\sigma \geq 528/125 \geq 4$ on the adverse recovery leg); the
index blindness of the composite reading at $w = (1, 1)$, whose minimum
stays at $2/5 > 0$ along the licensed plan while the ecological floor
reaches $-4/5 < 0$ mid-transition; the licensing thresholds
$\rho_1 = 2/3$ and $\rho_2 = 3/2$, cross-checked against the thresholds
derived directly from the tube geometry; and the rescue threshold
$\kappa^* = 1 - x$ on the non-typed-viable region, with shortfall $1/2$ at
the false-positive witness and full financing at the rescue witness
$x = 3/2$. All twenty-four checks pass in under a second on a laptop;
Table~\ref{tab:qa} collects the verification artifacts.

\begin{table}[t]
\centering
\small
\begin{tabular}{@{}ll@{}}
\toprule
Artifact & Result \\
\midrule
Benchmark exact checks (\texttt{safetransition verify}) & 24/24 pass \\
Unit and integration test suite & 27/27 pass \\
Runtime of checks + tests (laptop CPU) & under 1\,s each \\
Farkas certificate, stacked menu window & $\lambda = (1/2, 1/2)$, margin
$1/10$, re-verified \\
Belief recursion, recourse-failure system & root belief in
$\mathcal{W}_1 \setminus \mathcal{W}_2$ \\
Source / test lines & 1\,269 / 283 \\
\bottomrule
\end{tabular}
\caption{Verification artifacts shipped with SafeTransition 1.0.0. Every
check and certificate is evaluated in exact rational arithmetic.}
\label{tab:qa}
\end{table}

Figure~\ref{fig:benchmark} shows the benchmark panels regenerated from the
same verified values. Panel (a) is the separation a composite index cannot
see: the weighted reading stays certified over the full review period while
the floor it aggregates is driven below zero mid-transition. Panel (b) is
the plan menu as management schedules: the pulse-plus-closed-season quota,
the sustained-yield alternative at $\sigma(16/5) = 1088/125$, and the staged
rebuild below yield while the fund finances the buy-back.

\begin{figure}[t]
\centering
\includegraphics[width=\textwidth]{fig_benchmark.png}
\caption{The resource-transition benchmark, regenerated from the library's
exact checks. (a) The composite index at $w = (1,1)$ stays certified (min
$2/5 > 0$) along the licensed plan while the ecological margin $s_1$ is
breached mid-transition ($-4/5$ adverse, $-3/10$ benign): the separation the
typed operator exposes. (b) The plan menu as management schedules: FAST
(pulse $1463/125$ plus closed season), SLOW/NO-SWITCH (sustained yield
$1088/125$), and STAGED (below yield, rebuild to $69/20$), with the fund
schedule $x(t)$ under STAGED.}
\label{fig:benchmark}
\end{figure}

\subsection{Dashboard readings}\label{sec:readings}

The \texttt{compute} interface returns, at a state and weight, the per-plan
floor minima under each plan's characteristic disturbance, the composite
index along the read plan, the licensing verdicts at the weight ratio
$r = w_2/w_1$, the rescue threshold, and the index-blindness alarm, which
fires exactly when the composite minimum stays positive while a floor
minimum is negative. At the witness state $(x, s_1, s_2) = (1/2, 6/5, 6/5)$
with $w = (1,1)$, the alarm fires: the readings reproduce the benchmark
separation as a dashboard statement (Fig.~\ref{fig:readings}). The
\texttt{dashboard} module renders the readings, the benchmark status, and
any certificates as a single HTML file with inline SVG and CSS and no
external resources, so the file can be attached to a review, emailed, or
archived with a statutory record and remains readable offline. From the
reviewer's perspective this means the reading and its verification arrive
together: thresholds, alarms, and certificates in one file, each re-derivable
from the command line above.

\begin{figure}[t]
\centering
\includegraphics[width=\textwidth]{fig_readings.png}
\caption{Dashboard readings at the witness datum, exact. (a) Per-weight
licensing: the aggregate operator licenses FAST at weight ratios
$r \geq \rho_1 = 2/3$ and SLOW at $r \leq \rho_2 = 3/2$; the standard
equal-weight reading $r = 1$ sits inside the band where both are licensed
and no ratio licenses a typed-safe plan. (b) Rescue financing:
$\kappa^*(x) = \max(0,\, 1 - x)$; the false-positive witness $(1/2,\, 1/2)$
carries a shortfall of $1/2$, the rescue witness $(3/2,\, 0)$ finances the
buy-back with $1/2$ remaining.}
\label{fig:readings}
\end{figure}

\subsection{Partial observation and certificates}\label{sec:partial}

The bundled demonstrations exercise the certificate family on the companion
manuscript's instances. The three-state recourse failure runs the
belief-space recursion on a system whose root belief is one-step viable and
exits at the second step under every observation-based policy: the root
belief lies in $\mathcal{W}_1 \setminus \mathcal{W}_2$, witnessed only by
the recursion. The observation-fibre criterion reports the merged fibre on
which safe-set membership varies, ruling out any exact observation-only
certifier. The stacked menu window $u \leq 2/5$, $-u \leq -3/5$ returns the
Farkas certificate $\lambda = (1/2, 1/2)$ with infeasibility margin exactly
$1/10$, re-verified against $\lambda^{\top} A = 0$ and
$\lambda^{\top} b < 0$ before being returned; the certificate is the same
object the obstruction calculus attaches to the empty intersection of
compatible states' safe menus. The pooled-kernel identity of the deposited
double-integrator instance, weights $(3/8, 5/16, 5/16)$ with
$\sum_j \lambda_j n_j = 0$, is checked exactly.

\section{Relation to existing software}\label{sec:related}

SafeTransition complements three established families of tools, and borrows
from each. Numerical viability kernels and reachability tools (for the
viability-theoretic tradition, and for general-purpose simulation frameworks
in the environmental modelling literature, e.g.\ Raven
\citep{Craig2020} and FloPy \citep{Bakker2016}) compute set-valued safety
objects to high floating-point precision with impressive scalability; they
answer \emph{how large is the safe set}, where SafeTransition answers
\emph{is this verdict verifiable}, for the small, structured, rational
instances that assessment menus produce. Multi-criteria and sensitivity
analysis packages --- the SAFE toolbox \citep{Pianosi2015}, SALib
\citep{Herman2017}, SPOTPY \citep{Houska2015} --- share the library's
workflow-oriented design: model-independent interfaces, worked examples,
and benchmark functions; SPOTPY's demonstration that one package can carry
an entire family of methods with a minimal amount of glue code is the
presentation standard SafeTransition aims at. Modular environmental
simulators with public source code and regression suites --- Pywr
\citep{Tomlinson2020} and airGR \citep{Coron2017} --- inform the
availability discipline: permissive licensing, distribution through package
indexes, and a test suite executed on continuous integration. None of these
families returns certificates; none computes in exact arithmetic. The
comparison is functional, not competitive: the families operate at scales
and purposes for which floating-point computation is the right engineering
choice, and the exact layer becomes decisive precisely where verdicts enter
statutory or adversarial settings and must be re-derivable by third parties.

\section{Availability and requirements}\label{sec:availability}

\begin{table}[h]
\centering
\small
\begin{tabular}{@{}p{0.26\linewidth}p{0.68\linewidth}@{}}
\toprule
\textbf{Name} & SafeTransition \\
\textbf{Developer and contact} & Amin Abaee;
\href{mailto:amin_abaee@ut.ac.ir}{amin\_abaee@ut.ac.ir} \\
\textbf{Year first available} & 2026 \\
\textbf{Hardware required} & none beyond a standard computer \\
\textbf{Software required} & Python $\geq$ 3.9; no packages beyond the
standard library (matplotlib and Pillow only for regenerating figures;
publicly maintained third-party distributions are listed with their pinned
versions in the verification deposit) \\
\textbf{Program language} & Python \\
\textbf{Program size} & 1\,269 source lines, 283 test lines \\
\textbf{Availability and cost} & open source under the MIT licence; source,
tests, and examples archived in the verification deposit \citep{Deposit2026}
and mirrored in the project repository; no cost \\
\bottomrule
\end{tabular}
\caption{Software availability summary (EMS required fields).}
\label{tab:availability}
\end{table}

The test suite (27 tests) covers the benchmark, the operator chain and
witnessed separation on the witness datum, the certificate family on the
deposited instances, the belief-space recursion on the recourse-failure
system, and the dashboard's self-containment; it runs in under a second.
The dashboard example regenerates the single-file HTML from the library.

\section{Conclusions}\label{sec:conclusions}

SafeTransition turns the assessment-operator framework and the obstruction
calculus of the companion manuscripts into auditable software: typed
operators evaluated on exact tubes, recursions over state and belief space,
and infeasibility returned as re-verifiable Farkas certificates with exact
margins. The library's scope is deliberate: it certifies verdicts on small,
explicit, rational instances, and it makes no predictive claim --- the
theorem-level guarantees are those of the companion manuscripts, and on
continuous or belief-continuous systems the graph construction itself is
the modelling step. Within that scope, the tool closes a practical gap:
the readings that composite dashboards display, and the nonviability
verdicts that justify tightening a plan menu, can now carry their
certificates with them. Extensions under way are continuous-state data with
certified enclosures and additional certificate families from the
obstruction calculus; the interface intentionally leaves room for both.

\section*{CRediT authorship contribution statement}
\textbf{A.A.:} Conceptualization, Methodology, Software, Validation,
Formal analysis, Investigation, Data curation, Writing -- original draft,
Writing -- review \& editing, Visualization.

\section*{Declaration of generative AI in the writing process}
Deepseek AI and Qwen (Alibaba) assisted with code development. The author
reviewed and edited the content as needed and takes responsibility for the
content.

\section*{Acknowledgements}
Not applicable.

\bibliographystyle{elsarticle-harv}
\begin{thebibliography}{}
\providecommand{\natexlab}[1]{#1}

\bibitem[Abaee, 2026a]{Abaee2026a}
Abaee, A., 2026a. Aggregate Indices and Transition Safety: A Quantifier-Order
Separation Between Scalarized and Coordinate-Wise Feasibility. Manuscript
submitted for publication. Verified companion deposit:
\url{https://doi.org/10.6084/m9.figshare.33764023}.

\bibitem[Abaee, 2026b]{Abaee2026b}
Abaee, A., 2026b. An Obstruction Calculus for Viability under Incomplete
Observation. Manuscript submitted for publication.


\bibitem[Bakker et al., 2016]{Bakker2016}
Bakker, M., Post, V., Langevin, C.D., Hughes, J.D., White, J.T., Starn,
J.J., Fienen, M.N., 2016. Scripting MODFLOW model development using open
source Python and FloPy. Groundwater 54 (5), 733--739.
\url{https://doi.org/10.1111/gwat.12413}.

\bibitem[Coron et al., 2017]{Coron2017}
Coron, L., Thirel, G., Delaigue, O., Perrin, C., Andr\'eassian, V., 2017.
The suite of lumped GR hydrological models in an R package. Environmental
Modelling \& Software 94, 166--171.
\url{https://doi.org/10.1016/j.envsoft.2017.05.002}.

\bibitem[Craig et al., 2020]{Craig2020}
Craig, J.R., Brown, G., Chlumsky, R., Jenkinson, R.W., Jost, G., Lee, K.,
 Mai, J., Serrer, M., Sgro, N., Shafii, M., Snowdon, A.P., Tolson, B.A.,
 2020. Flexible watershed simulation with the Raven hydrological modelling
 framework. Environmental Modelling \& Software 129, 104728.
 \url{https://doi.org/10.1016/j.envsoft.2020.104728}.

\bibitem[Abaee, 2026]{Deposit2026}
Abaee, A., 2026. Verification code and figure pipeline for Aggregate
Indices and Transition Safety [data set]. figshare.
\url{https://doi.org/10.6084/m9.figshare.33764023}.

\bibitem[Herman and Usher, 2017]{Herman2017}
Herman, J., Usher, W., 2017. SALib: An open-source Python library for
sensitivity analysis. Journal of Open Source Software 2 (9), 97.
\url{https://doi.org/10.21105/joss.00097}.

\bibitem[Houska et al., 2015]{Houska2015}
Houska, T., Kraft, P., Chamorro-Chavez, A., Breuer, L., 2015. SPOTting
model parameters using a ready-made Python package. PLOS ONE 10 (12),
e0145180. \url{https://doi.org/10.1371/journal.pone.0145180}.

\bibitem[Jakeman et al., 2006]{Jakeman2006}
Jakeman, A.J., Letcher, R.A., Norton, J.P., 2006. Ten iterative steps in
development and evaluation of environmental models. Environmental Modelling
\& Software 21 (5), 602--614. \url{https://doi.org/10.1016/j.envsoft.2006.01.004}.

\bibitem[Pianosi et al., 2015]{Pianosi2015}
Pianosi, F., Sarrazin, F., Wagener, T., 2015. A Matlab toolbox for global
sensitivity analysis. Environmental Modelling \& Software 70, 80--85.
\url{https://doi.org/10.1016/j.envsoft.2015.04.009}.

\bibitem[Tomlinson et al., 2020]{Tomlinson2020}
Tomlinson, J.E., Arnott, J.H., Harou, J.J., 2020. A water resource
simulator in Python. Environmental Modelling \& Software 126, 104635.
\url{https://doi.org/10.1016/j.envsoft.2020.104635}.

\end{thebibliography}

\end{document}
